\section{Why AR's Compactness Is Legally Defensible}
\label{sec:legal}

\subsection{The Three-Part Test}

We propose a three-part test for the legal defensibility of an algorithmic
plan's compactness position:

\begin{enumerate}
  \item \textbf{Above the ensemble mean}: The plan should be more compact
    than the average valid plan, demonstrating that compactness was a genuine
    objective.
  \item \textbf{Below the 95th percentile}: The plan should not be at the
    extreme of the compactness distribution, which would invite scrutiny
    for partisan packing.
  \item \textbf{Explainable objective}: The source of the compactness should
    be traceable to a stated objective (edge-cut minimisation), not to a
    post-hoc compactness filter applied after partisan optimisation.
\end{enumerate}

The \ar{} plan satisfies all three conditions in all studied states.

\subsection{Applying the Test}

\paragraph{Condition 1: Above the ensemble mean.}
The \ar{} plan's $\bar\pp$ exceeds the ensemble mean in all six states
by margins ranging from $0.065$ (California) to $0.087$ (North Carolina).
This margin is statistically significant ($p < 0.001$) given the ensemble
standard deviation of approximately $0.031$.

\paragraph{Condition 2: Below the 95th percentile.}
As shown in Section~\ref{sec:outlier}, no \ar{} state plan exceeds the
ensemble 95th percentile.  The plan is compact but not an outlier.

\paragraph{Condition 3: Explainable objective.}
The METIS edge-cut objective is publicly documented and mathematically
defined \citep{karypis1998fast}.  It is applied before any partisan data
is consulted.  The compactness of the resulting plans is a consequence
of the objective, not a post-hoc selection criterion.

\subsection{Contrast with Human-Drawn Maps}

Enacted congressional maps often fail one or more conditions of this test.
The North Carolina plan evaluated by \citet{herschlag2020quantifying} fell
at the 99th percentile of Republican seats (failing condition 2 on partisan
grounds, not compactness) and also had below-median compactness (failing
condition 1).

The \ar{} plan presents the opposite profile: it is above median on
compactness but near median on partisan outcomes.  This combination is
exactly what an independent redistricting commission would seek to achieve,
and what courts have cited as evidence of good-faith map-drawing.

\subsection{The \texorpdfstring{\shaw{}}{Shaw} Standard}

\shaw{} held that districts of ``bizarre shape'' that could not be explained
by traditional redistricting criteria were subject to strict scrutiny.  The
\ar{} plan's position at the $61$st--$72$nd percentile of compactness
provides strong evidence that the districts are not bizarre in the \shaw{}
sense: they are more compact than average, meaning their shapes are well
within the range of what geographic criteria produce.

Importantly, \shaw{} was about racial gerrymanders, and the \ar{} algorithm
does not use racial data.  But the compactness evidence is useful as
corroborating evidence that the plan's boundaries follow geographic rather
than demographic logic.
